\section{代表论文介绍与评析}

\writingguide{建议写 1500--1800 字。这是作业要求中的“某篇论文介绍”。需阅读论文原文，而不是只依据摘要。建议覆盖研究动机、方法、数据集、实验设置、结果、贡献、局限和个人评价。}

本节暂以 WWW 2026 论文 \emph{CFVBench: A Comprehensive Video Benchmark for Fine-grained Multimodal Retrieval-Augmented Generation} 为示例\parencite{wei2026cfvbench}。如果最终选择其他论文，应同步修改本节标题内容、图表和参考文献。

\subsection{研究问题与动机}

\placeholder{用自己的话说明论文试图解决的问题、已有工作的不足，以及该问题与本 Track 的关系。}

\subsection{核心方法}

\placeholder{解释整体流程、模型结构或数据集构建方法。先给出直观解释，再介绍必要的技术细节。}

\ifdraftmode
\begin{figure}[htbp]
  \centering
  \fbox{\parbox[c][4.2cm][c]{0.82\textwidth}{\centering
    在此放置自己绘制的方法框架图，或使用\\
    \texttt{\textbackslash includegraphics[width=...]} 插入图片。}}
  \caption{代表论文的方法框架（占位图，提交前替换）}
  \label{fig:method-overview}
\end{figure}
\fi

\subsection{实验设计与主要结果}

\placeholder{介绍数据集、对比方法、评价指标和主要实验结果。不要只说“效果更好”，应给出关键数值并解释其意义。}

\begin{table}[htbp]
  \centering
  \caption{代表论文关键实验结果整理（示例结构）}
  \label{tab:main-results}
  \begin{tabular}{lccc}
    \toprule
    方法 & 指标一 & 指标二 & 指标三 \\
    \midrule
    基线方法 & -- & -- & -- \\
    论文方法 & -- & -- & -- \\
    \bottomrule
  \end{tabular}
\end{table}

\subsection{贡献、局限与个人评价}

\placeholder{区分作者声称的贡献与自己的评价；结合实验设计分析结论是否充分，并指出适用边界、潜在偏差和改进方向。}
